%
% ###################################################################
%  Copyright (c) 2026, REPLACE: Your Name
%  Editors: A.D. Rollett & M. De Graef
%  All rights reserved.
%
% Licensed under the Creative Commons CC BY-NC-SA 4.0 License,
% hereafter referred to as the "License"; you may not use this
% document except in compliance with the License. You may obtain
% a copy of the License at
%     https://creativecommons.org/licenses/by-nc-sa/4.0/legalcode
% Unless required by applicable law or agreed to in writing, all
% material distributed under the License is distributed on an
% "AS IS" BASIS, WITHOUT WARRANTIES OR CONDITIONS OF ANY KIND,
% either express or implied. See the License for the specific
% language governing permissions and limitations under the License.
% ###################################################################

% ###################################################################
% AUTHOR SETUP — Edit the lines below
% ###################################################################

% Uncomment the next line ONLY if this is a foundational chapter:
%\corechapter{Yes}

% Your name and affiliation (appears below the chapter title):
\OCchapterauthor{REPLACE: Your Name, Your University}

% Chapter abbreviation: choose a UNIQUE 6-character uppercase code.
% This is used as a prefix for all labels in this chapter.
% Examples: LINALG, TENSOR, NUMSYS, PROJTS
\renewcommand{\chabbr}{MYCHAP}

% Chapter header image (2480 x 1240 pixels, 300 dpi, RGB format).
% Replace \noheaderimage with your header filename (e.g., MYCHAPheader.pdf).
% The file should be placed in chapter/pdf/.
% Note that \graphicspath is defined in the main.tex file.
%\chapterimage{\graphicspath MYCHAPheader.pdf}
\chapterimage{\noheaderimage}

% ###################################################################
% CHAPTER TITLE — Edit the title below
% ###################################################################
\chapter{My Chapter Title}\label{\chabbr:ch:MyChapter}

% Register author information (repeat for each co-author)
\writeauthor{\chabbr:ch:MyChapter}{My Chapter Title}{REPLACE: Last}{REPLACE: First}{REPLACE: Department}{REPLACE: University}{REPLACE: email@example.com}{REPLACE: https://homepage.example.com}

% ###################################################################
% LEARNING OBJECTIVES
% ###################################################################
% Each chapter begins with Learning Objectives.  Each objective should:
% - Contain an active verb (Learn, Understand, Derive, Apply, etc.)
% - Link to the relevant section using the \chabbr:sec:label format
% - Use the OCBurntOrange color for the reference
\begin{learningobjectives}
{\begin{itemize}
    \item[{\color{OCBurntOrange}\ref{\chabbr:sec:introduction}:}] Understand the basic concepts presented in this chapter
    \item[{\color{OCBurntOrange}\ref{\chabbr:sec:methods}:}] Learn the methods used for analysis
    \item[{\color{OCBurntOrange}\ref{\chabbr:sec:examples}:}] Apply the concepts to worked examples
\end{itemize}}
\end{learningobjectives}

% ###################################################################
% CHAPTER CONTENT STARTS HERE
% ###################################################################

\section{Introduction}\label{\chabbr:sec:introduction}

This is the introduction to your chapter.  You can use standard \LaTeX\ commands
for formatting, equations, figures, and tables.  All labels in this chapter should
use the prefix defined by \texttt{\textbackslash chabbr}, which is \texttt{\chabbr}
for this template.

The label format is \texttt{\textbackslash chabbr:TYPE:name}, where TYPE is one of:
\texttt{ch} (chapter), \texttt{sec} (section), \texttt{ssec} (subsection),
\texttt{eq} (equation), \texttt{fig} (figure), \texttt{tb} (table).

\subsection{An Example Subsection}\label{\chabbr:ssec:example}

Here is an example of an inline equation: $E = mc^2$.  And here is a displayed
equation with a label:
\begin{equation}
    \mathbf{F} = m\mathbf{a}
    \label{\chabbr:eq:newton}
\end{equation}
We can refer to Eq.~\ref{\chabbr:eq:newton} anywhere in this chapter.

% ###################################################################
\section{Methods}\label{\chabbr:sec:methods}

This section describes the methods.  Here is an example of a figure reference.
Figure~\ref{\chabbr:fig:example} shows an example plot from the number systems
chapter.

\begin{figure}[htbp]
    \centering
    \includegraphics[width=0.6\textwidth]{\graphicspath figcomplex.pdf}
    \caption{\label{\chabbr:fig:example}An example figure showing the complex
    plane.  Replace this with your own figure.  Place PDF figures in the
    \texttt{chapter/pdf/} directory.}
\end{figure}

Here is an example of a table:
\begin{table}[htbp]
    \centering
    \caption{\label{\chabbr:tb:example}An example table with three columns.}
    \begin{tabular}{lcc}
        \toprule
        Property & Symbol & Units \\
        \midrule
        Force & $F$ & N \\
        Mass & $m$ & kg \\
        Acceleration & $a$ & m/s$^2$ \\
        \bottomrule
    \end{tabular}
\end{table}

If you want to add an word to the index, you can use the standard \texttt{\textbackslash index} command with that word as argument; alternatively,
use the \texttt{\textbackslash indexit} command to have the word in italics in the text and add it to the index with a single command. Let's add \indexit{index} to the index.

% ###################################################################
\section{Worked Examples}\label{\chabbr:sec:examples}

This section contains worked examples.  You can use the
\texttt{OCexample} environment for formatted examples:

 If you want to cite a reference, use \texttt{\textbackslash cite} with a key from
 chapter/chaptercitations.bib.  Example: \cite{author2026a}.

% ###################################################################
% END OF CHAPTER
% ###################################################################
